\begin{pastexam}[label={px:ncku-115-prog-16}]{成大資訊聯招 115 學年度 演算法問答第 3 題}
（問答）

[10\%] You are analyzing the routing latencies for a constellation of 4 satellites (Nodes A, B, C, D).
Below is the adjacency list with the latency between specific nodes:

\textbf{Adjacency List:}
\begin{itemize}[label=$\blacksquare$,nosep]
\item A $\to$ B: 10
\item A $\to$ C: -5
\item C $\to$ B: 4
\item B $\to$ D: 5
\item D $\to$ C: 2
\item D $\to$ A: 3
\end{itemize}

\textbf{Requirement:} You need to analyze the latency characteristics of the network so that shortest
communication paths can be efficiently determined when needed

\underline{\textbf{Questions}}:
\begin{enumerate}[label=(\arabic*)]
\item {[4\%]} Based on the specific properties of the edges in this graph and the problem requirement,
identify the \textbf{single most appropriate standard algorithm} to use. \textbf{Explain your choice}
by pointing out the specific graph feature that disqualifies other algorithms, and explicitly state
\textbf{why} at least one other commonly used shortest-path algorithm would fail or be inappropriate in
this case.
\begin{itemize}
\item Algorithm: \underline{\hspace{6cm}}
\item Why choosing this algorithm? \underline{\hspace{4.5cm}}
\end{itemize}
\item {[2\%]} Let $M$ be the initial $\mathbf{4\times 4}$ \textbf{distance matrix} used at the start of
your chosen algorithm. Please construct and write down $M$ clearly.\\
(Note: Use ``$\infty$'' for non-existing edges and ``0'' for self-loops. Order: A, B, C, D)
\begin{itemize}\item $M = \left[\rule{0pt}{2.2em}\hspace{4cm}\right]$\end{itemize}
\item {[4\%]} After the algorithm completes execution
\begin{enumerate}[label=(\alph*)]
\item {[2\%]} What is the \textbf{final shortest path cost} from Node D to Node B? Please
\textbf{trace the path logic} (e.g., D $\to$ \dots $\to$ B).
\begin{itemize}
\item Final shortest path cost: \underline{\hspace{6cm}}
\item Final path logic: \underline{\hspace{6cm}}
\end{itemize}
\item {[2\%]} State the asymptotic \textbf{Time Complexity} $\Theta$ of this algorithm.
\begin{itemize}\item Time Complexity: \underline{\hspace{6cm}}\end{itemize}
\end{enumerate}
\end{enumerate}
\end{pastexam}

\begin{pastsoln}
答案：(1) Floyd--Warshall；(2) 見下；(3)(a) $2$，D $\to$ A $\to$ C $\to$ B；(b) $\Theta(V^3)$，這裡
$V = 4$。

\textbf{(1) 選 Floyd--Warshall。} 圖有兩個特徵決定了選擇：
\begin{itemize}
\item \textbf{有負權重的邊} A $\to$ C $= -5$。 Dijkstra 演算法假設邊權重非負：一個點被取出後距離就
定案，負邊可能在之後讓它變得更短，Dijkstra 卻不會再更新（資結書第~25.2.5~節）。 實際在這張圖上從 D
跑 Dijkstra：先定案 D $= 0$，鬆弛得 C $= 2$、A $= 3$；取出 C 定案為 $2$，鬆弛得 B $= 6$；取出 A，
A $\to$ C 給出 $3 - 5 = -2 < 2$，但 C 已經定案；最後 B 定案為 $6$。 真正的 D 到 B 是 $2$，Dijkstra 錯了。
\item \textbf{沒有負環}，所以最短路徑有定義。 檢查所有的環：A $\to$ C $\to$ B $\to$ D $\to$ A 是
$-5 + 4 + 5 + 3 = 7$，A $\to$ B $\to$ D $\to$ A 是 $18$，C $\to$ B $\to$ D $\to$ C 是 $11$，都是正的。
\end{itemize}
題目要的是「需要時能快速查到任兩點的最短路徑」，也就是所有點對最短路徑，算好一張表之後每次查詢
$O(1)$。 能處理負邊的所有點對演算法裡，Floyd--Warshall 最直接（資結書第~25.6~節），$V^3 = 64$ 次運算。
Bellman--Ford 也能處理負邊，但它是單一起點的，要所有點對得跑 $V$ 次，$O(V^2 E)$，在稠密圖上是
$O(V^4)$；Johnson 演算法適合稀疏的大圖，對 $4$ 個點太複雜。

\textbf{(2) 初始矩陣}（列是起點、行是終點，順序 A, B, C, D）：
\[
M = D^{(0)} =
\begin{bmatrix}
0 & 10 & -5 & \infty \\
\infty & 0 & \infty & 5 \\
\infty & 4 & 0 & \infty \\
3 & \infty & 2 & 0
\end{bmatrix}.
\]

\textbf{(3)(a)} Floyd--Warshall 依序讓 A、B、C、D 可以當中間點，
\[
D^{(k)}[i][j] = \min\bigl\{D^{(k-1)}[i][j],\; D^{(k-1)}[i][k] + D^{(k-1)}[k][j]\bigr\}
\]
（資結書虛擬碼~25.8）。 粗體是這一輪變小的格子：
\[
D^{(1)}_{\text{經 A}} = \begin{bmatrix}
0 & 10 & -5 & \infty \\
\infty & 0 & \infty & 5 \\
\infty & 4 & 0 & \infty \\
3 & \mathbf{13} & \mathbf{-2} & 0
\end{bmatrix}\!,\quad
D^{(2)}_{\text{經 B}} = \begin{bmatrix}
0 & 10 & -5 & \mathbf{15} \\
\infty & 0 & \infty & 5 \\
\infty & 4 & 0 & \mathbf{9} \\
3 & 13 & -2 & 0
\end{bmatrix}\!,
\]
\[
D^{(3)}_{\text{經 C}} = \begin{bmatrix}
0 & \mathbf{-1} & -5 & \mathbf{4} \\
\infty & 0 & \infty & 5 \\
\infty & 4 & 0 & 9 \\
3 & \mathbf{2} & -2 & 0
\end{bmatrix}\!,\quad
D^{(4)}_{\text{經 D}} = \begin{bmatrix}
0 & -1 & -5 & 4 \\
\mathbf{8} & 0 & \mathbf{3} & 5 \\
\mathbf{12} & 4 & 0 & 9 \\
3 & 2 & -2 & 0
\end{bmatrix}\!.
\]
D 到 B 的變化：一開始沒有直接的邊（$\infty$）；經 A 時 D $\to$ A $\to$ B $= 3 + 10 = 13$；經 C 時
D $\to$ C 已經是 $-2$（第一輪得到的 D $\to$ A $\to$ C），再接 C $\to$ B：$-2 + 4 = 2 < 13$。 最後一輪
沒有再變小。 所以最短距離是 $2$，路徑
\[
\text{D} \to \text{A} \to \text{C} \to \text{B}: \quad 3 + (-5) + 4 = 2 .
\]
直接走 D $\to$ C $\to$ B 是 $2 + 4 = 6$，繞到 A 去走那條 $-5$ 的邊反而更短。 對角線最後都是 $0$，也確認了
沒有負環。

\textbf{(3)(b)} 三層迴圈，$k$、$i$、$j$ 各跑 $V$ 次，每次 $O(1)$：$\Theta(V^3)$，空間 $\Theta(V^2)$。
\end{pastsoln}
